% =====================================================================
% Trabalhos relacionados --- Vitória (Frente 2)
% Só a Vitória edita este arquivo.
% =====================================================================
\section{Trabalhos relacionados}\label{sec:trabalhos}

Simulando epidemias SIR em redes reais, um estudo clássico mostrou que o número de núcleo
($k$-core) prevê os melhores espalhadores melhor do que o grau e a intermediação: o que importa é
estar no miolo denso da rede, e não só ter muitos vizinhos~\cite{kitsak2010}. Uma revisão posterior
reuniu dezenas de métodos para encontrar nós importantes, do grau ao PageRank e a métodos de
aprendizado, e os comparou pelo alcance em simulações de espalhamento; nenhuma medida venceu em
todas as redes~\cite{lu2016}. No BraSNAM, grafos de retweets dos protestos brasileiros de 2013
mostraram que o PageRank e a centralidade de autovetor destacam usuários que a simples contagem de
retweets não revela~\cite{theodoro2014}, e, numa rede de pilotos de Fórmula 1, o PageRank deu o
ranking mais informativo entre as centralidades comparadas~\cite{faria2024}.

Este trabalho segue a pergunta do estudo de espalhadores~\cite{kitsak2010}, mas numa rede brasileira, a Twitch em português,
e com dois gabaritos em vez de um: o alcance no SIR e as visualizações reais de cada canal. Além
disso, as seis centralidades comparadas e a simulação foram implementadas do zero, em Python puro,
e conferidas contra valores publicados e contra o NetworkX~\cite{hagberg2008}.